\section{Ordered adaptive integration with rightmost-root tracking}\label{sec:method}

\subsection{The winding integral as an ordinary differential equation}\label{sec:phaseode}

The key step is disarmingly simple: instead of applying quadrature
to~\eqref{eq:count}, the accumulated phase is treated as the state of the
scalar initial value problem
\begin{equation}\label{eq:phaseode}
  y'(\omega) = \theta'(\omega)
             = \Imag\frac{\Dfun'(\omega)}{\Dfun(\sigma+\ii\omega)},
  \qquad y(0)=0,
  \qquad
  \Zraw = \frac{n}{2}-\frac{y(\wmax)}{\pi},
  \qquad
  \Zint = \operatorname{round}(\Zraw),
\end{equation}
which is handed to a modern adaptive integrator with local error control.
It should be stated plainly what this formulation is and is not. The
right-hand side of~\eqref{eq:phaseode} does not depend on the state $y$, so
the Jacobian $\partial_y y'$ vanishes identically: \eqref{eq:phaseode} is a
quadrature problem in disguise, it is never stiff in the classical sense,
and an automatically stiffness-switching integrator provably never switches
on it. What the integrator contributes is therefore not stiffness handling
but \emph{adaptive step-size control on an ordered march}: steps contract
around the near-singular peaks~\eqref{eq:peak}, where the phase turns by
$\pm\pi$ over a frequency interval of width $|\sigma_r-\sigma|$, and expand
again over the smooth stretches. Because the integrand is analytic, a
high-order explicit Runge--Kutta pair is the appropriate choice; the default
is \code{Vern9}~\cite{verner2010numerically}, which at the default tolerance
is three orders of magnitude more accurate than the fifth-order
alternative~\cite{tsitouras2011runge} for a factor of $2.4$ in cost
(Section~\ref{sec:performance-convergence}). The ordered march is
what makes the root tracking of Section~\ref{sec:tracking} possible, which
is the reason the phase ODE is preferred here over plain quadrature
(Section~\ref{sec:alternatives}).

The integrand is evaluated \emph{exactly}: the characteristic function is called
once per frequency with a dual-number argument
$\lam = \sigma + \ii(\omega+\delta)$, $\delta^2=0$, so
forward-mode automatic differentiation~\cite{revels2016forward} delivers
$\Dfun$ and $\Dfun'=\partial_\omega\Dfun$ simultaneously, without finite
differencing and without any user-supplied derivative. A tiny offset
$\omega\mapsto\max(\omega,10^{-9})$ regularizes branch points of fractional
terms at the origin.

Two practical parameters remain. The truncation frequency $\wmax$ can be
chosen generously: past the highest structural resonance the phase slope
decays algebraically, the adaptive step size grows accordingly, and a
\emph{too small} $\wmax$ is immediately visible in the result, because the
truncated tail shifts $\Zraw$ away from an integer
(Section~\ref{sec:integerness}). The cost of the tail is governed by the
decay of the delay-induced phase ripple. If the delayed terms sit at least
two orders below the leading power of $\Dfun$ -- as for delayed
\emph{position} feedback on mechanical systems -- the ripple decays at least
like $\omega^{-2}$ and pushing $\wmax$ far beyond the resonances is
essentially free: for the fourth-order benchmark the march costs $1426$,
$1474$ and $1506$ evaluations of $\Dfun$ at $\wmax=10^{6}$, $10^{8}$ and
$10^{10}$ respectively -- four extra decades for six percent more work. If a
delayed term sits only one order below the leading power (delayed
\emph{velocity} feedback), the ripple decays merely like $\omega^{-1}$, the
solver keeps resolving it (Fig.~\ref{fig:walkthrough}c), and $\wmax$ does
cost. The resulting recommendation is asymmetric: \emph{for accuracy, be
generous} ($\wmax=10^{6}$ is a safe default for every system in this paper);
\emph{for speed, be frugal} -- two to three decades above the highest
resonance leaves a truncated tail of order $0.2/\wmax$ in $\Zraw$, far below
the rounding threshold, and for a slowly decaying ripple cuts the evaluation
count by more than an order of magnitude. Since a too-small $\wmax$
announces itself in the residual, the choice can be made empirically rather
than by analysis.

The leading order $n$ in~\eqref{eq:phaseode} is estimated numerically, and
the natural place to do so is not the imaginary axis at all. The term $n/2$
accounts for the phase gathered on the arc at infinity, and that arc lies in
the \emph{right half-plane}, where every delayed factor $\ee^{-\lam\tau_j}$
decays exponentially. Probing along the positive real axis therefore
measures precisely the quantity the counting formula needs, with the delay
terms already extinguished:
\begin{equation}\label{eq:npow}
  n \;\approx\; \operatorname*{LS\text{-}slope}_{\;s_k\in[s_\ast/10,\;s_\ast]}
  \Bigl(\ln s_k \,\mapsto\, \ln\bigl|\Dfun(s_k)\bigr|\Bigr),
  \qquad s_\ast \gg 1 ,
\end{equation}
the least-squares slope of the log--log magnitude over a handful of
log-spaced real samples. The probe is decoupled from the integration range
($n$ is a property of $\Dfun$ alone; the default is $s_\ast=10^{8}$ whatever
$\wmax$ is), costs a handful of evaluations once per sweep, and converges as
$O(1/s_\ast)$ -- measured here as $7.5\cdot10^{-5}$, $7.4\cdot10^{-7}$ and
$7.4\cdot10^{-9}$ at $s_\ast=10^{4}$, $10^{6}$ and $10^{8}$ -- including the
effective integer order of neutral equations and the non-integer order of
fractional ones. The real axis is essential here: on the imaginary axis
$|\ee^{-\ii\omega\tau}|=1$, so the delay terms superimpose a ripple on
$\ln|\Dfun(\ii\omega)|$ that a fit aliases rather than averages, and for a
neutral system, whose ripple never decays, a single-point estimator
diverges in proportion to $\omega$ (measured: $n_{\mathrm{est}}=80\,028$ at
$\omega=10^{6}$ for a system whose true effective order is $2$). Probing
where the delays are exponentially small removes all of this at a stroke,
and in parameter sweeps the estimate is computed once when $n$ is
parameter-independent.

Two limitations are worth naming. First, $|\Dfun(s)|\sim s^{n}$ overflows
for large $n$, so $s_\ast$ is reduced automatically until $|\Dfun|$ is
finite, which for very high-order determinants caps the achievable accuracy.
The better answer is structural rather than numerical: whenever the model
offers a natural open-loop/feedback split, the \emph{return difference}
$\Dfun=\det(\mat I+\mat Q_0^{-1}\mat F)=1+L(\lam)$ has the same
right-half-plane zeros, adds poles only at the (stable) open-loop roots, and
tends to $1$ at infinity -- then $n=0$ \emph{exactly} and the enormous
magnitudes disappear; case study A.9 in Appendix~\ref{app:gallery}
quantifies the difference. Second, Equation~\eqref{eq:npow} presumes that
$\Dfun$ is polynomially bounded in the right half-plane, which holds for
every quasi-polynomial but not for the transcendental characteristic
functions of continua, where $\cosh\bigl(\sqrt{\lam^2+c\lam}\,L\bigr)$ grows
like $\ee^{\lam L}$. The package detects this automatically (for a genuine
power law the fitted slope is reproduced one window lower; for exponential
growth it is not) and falls back to a fit along the imaginary axis, where
such functions are bounded ($n\approx0$ for the beam of
Appendix~\ref{app:gallery}).

The requested relative and absolute solver tolerances control the local
error of the accumulated phase, measured in radians. Since the exact
winding number is an integer multiple of $\pi$ in the phase, the error of
$\Zraw$ inherits the tolerance directly as a \emph{dimensionless} quantity:
the same default ($10^{-5}$) is appropriate for a machining model and for a
dense $50\times50$ matrix determinant alike, with no problem-dependent
scaling (Section~\ref{sec:integerness}).

\subsection{Rightmost-root estimation as a by-product}\label{sec:tracking}

The peaks of the integrand are not only a numerical nuisance -- they carry
the information where the dominant characteristic roots are. Along the
integration line, $|\Dfun(\sigma+\ii\omega)|^{2}$ dips to a sharp local
minimum at $\omega\approx\omega_r$ for every nearby root
$\lam_r=\sigma_r+\ii\omega_r$, and the smallest minimum belongs to the root
closest to the line, which for $\sigma=0$ is typically the rightmost root
that decides stability. The solver therefore tracks
$\min_\omega|\Dfun|^{2}$ during the march at negligible cost: at every
integrand evaluation the already computed quantities $\Dfun$ and $\Dfun'$
yield the real part of one Newton step towards the root,
\begin{equation}\label{eq:sigmaest}
  \sest \;=\; \sigma + \Real\!\left(-\frac{\Dfun}{\dd\Dfun/\dd\lam}\right)
        \;=\; \sigma - \frac{\Real\Dfun\cdot\partial_\omega\Imag\Dfun
              - \Imag\Dfun\cdot\partial_\omega\Real\Dfun}{|\Dfun'|^{2}},
\end{equation}
using $\dd\Dfun/\dd\lam=-\ii\,\Dfun'$, and the frequency of the minimum is
stored as $\west$. The pair $(\sest,\west)$ is a first-order estimate of the
dominant root location returned together with $\Zint$; its accuracy is
highest exactly at the stability boundary, where the root lies on the line
and the Newton step becomes exact up to the solver tolerance. Away from the
boundary the estimate degrades gracefully, but there only its sign and
magnitude matter -- as a smooth measure of the distance to instability.

Several roots can be tracked simultaneously. A compile-time variant of the
solver detects every local minimum of $|\Dfun|^{2}$ through the sign change
of its $\omega$-derivative, $2(\Real\Dfun\,\partial_\omega\Real\Dfun +
\Imag\Dfun\,\partial_\omega\Imag\Dfun)$, and maintains the $N_r$ best
candidates in a sorted, allocation-free buffer; a boundary test at
$\omega=0$ captures real dominant roots, since $|\Dfun|^{2}$ is even in
$\omega$. Tracking one root is the default and adds a few percent to the sweep
time; tracking
zero roots gives maximal sweep speed. One caveat deserves emphasis: the
\emph{deepest} minimum of $|\Dfun|$ belongs to the rightmost root near the
stability boundary, but deeper inside the stable or unstable domains a
different root branch may produce a deeper dip. When the true dominant root
is required away from the boundary, several minima should be tracked
($N_r\approx5$--$10$), all of them refined, and the dominant root taken as
the one with the maximal real part -- this recipe is validated against an
independent eigenvalue method in Section~\ref{sec:performance-sd}.

That recommendation is worth more than a footnote, because the single-minimum
estimate is not merely inaccurate away from the boundary: it is
\emph{discontinuous}. Where one branch overtakes another as the closest one,
$\west$ jumps between modes and $\sest$ jumps with it, and a method that fits
a local linear model -- as the boundary tracer of Section~\ref{sec:mdbm} does
inside each cell -- will read the jump as a zero crossing and place a
spurious boundary point in the middle of the stable domain. The dominant root
passes through a branch crossing smoothly, because it is a maximum over
branches rather than the nearest one, so building the coloring and the
boundary objective from it removes the artefact at its source. For the count
itself the caveat is immaterial, since near the boundary the crossing root
always produces the global minimum.

\subsection{Optional higher-order root refinement}\label{sec:refinement}

The raw tracked estimate $(\sest,\west)$ -- one Newton step off the minimum
of $|\Dfun|$ -- is the method's \emph{default} root output, and deliberately
so. It is first-order accurate, but it is the \emph{smoothest} field over a
parameter chart (a single formula evaluated at every pixel), it is the
cheapest, and its \emph{sign} -- all the stability classification needs -- is
already correct. When an accurate rightmost root is wanted, the same automatic
derivatives drive an optional local polish. Writing the Taylor expansion at
$\hat\lam=\sest+\ii\west$,
\begin{equation}\label{eq:taylor}
  \Dfun(\hat\lam+\Delta) \approx \Dfun + \Dfun_\lam\,\Delta
   + \tfrac12 \Dfun_{\lam\lam}\,\Delta^{2}
   + \tfrac16 \Dfun_{\lam\lam\lam}\,\Delta^{3} = 0,
\end{equation}
where the subscripts denote derivatives with respect to $\lam$ evaluated at
$\hat\lam$, the package offers (i) \emph{Newton--Raphson} iterations
$\Delta=-\Dfun/\Dfun_\lam$; (ii) a single \emph{quadratic} (Poly-2) or (iii)
\emph{cubic} (Poly-3) Taylor solve, where the root of the local polynomial
closest to $\hat\lam$ is selected via the eigenvalues of its companion matrix;
and (iv) a \emph{combined} mode that runs the polynomial and Newton polishes
and keeps, per root, the candidate with the smallest $|\Dfun|$. The first
derivative is exact (automatic differentiation); the higher derivatives are
central differences of that exact first derivative, which avoids the
compilation cost of nested dual numbers. Each polish is \emph{monotone}: it is
accepted only if it lowers $|\Dfun|$ below the seed, otherwise the raw estimate
is returned, so refinement can help but never degrade the estimate.

Which polish is best is genuinely system-dependent: on a smooth
quasi-polynomial the Newton iteration reaches machine precision, but on a
rational characteristic function, near the poles of the plant frequency
response, a Newton step can settle on a neighbouring root and speckle an
otherwise smooth chart, where the Taylor polishes stay put. The combined
mode is the robust choice -- lowest residual on every system tested -- at
about $\RefCombOver\%$ over the raw solve; the trade-offs are measured in
Section~\ref{sec:showcase-roots}. The practical guidance matches the design
goal of a single, generally valid setting: leave refinement off for charts,
turn it on when a specific rightmost root matters, and read a precise
decay-rate \emph{map} from the $\sigma$-level contours of
Section~\ref{sec:sigmacontours} rather than from a per-pixel polish.

If several estimates are tracked, the same refinement applied to each returns
a set of dominant roots, comparable to the output of dedicated root finders
such as QPmR~\cite{vyhlidal2009mapping} -- with the difference that here the
root solve is a cheap, optional post-processing step of the stability count.

\subsection{Self-validating error estimate}\label{sec:integerness}

The exact winding number is an integer, so
\begin{equation}\label{eq:eps}
  \varepsilon = \bigl|\Zraw - \operatorname{round}(\Zraw)\bigr|
\end{equation}
is a rigorous, freely available indicator of the accumulated numerical
error: it contains the local integration errors, the tail truncated at
$\wmax$, and the residual of the estimated order~\eqref{eq:npow}. (When a
characteristic root lies on the integration line itself, no count is defined;
the package then returns an invalid-count marker rather than a number.)
Since the
phase is integrated in radians, $\varepsilon$ is \emph{dimensionless} and
scales essentially linearly with the requested tolerance -- amplified by
the accepted-step count, since the per-step error tolerance accumulates
along the march -- hence the tolerance acts as a universal accuracy dial
that transfers between systems without re-tuning. In the numerical
experiments of Section~\ref{sec:showcase-tolerance}, $\varepsilon$ tracks
the tolerance across six decades: at $10^{-2}$ the residual reaches order
one and genuine miscounts appear -- the deliberate failure demonstration of
Fig.~\ref{fig:errorfield} -- while $10^{-4}$ already keeps $\varepsilon$ an
order of magnitude below the rounding threshold $1/2$ and produces reliable
charts at near-maximal speed, and $10^{-8}$ yields reference-quality values
at a moderate CPU-time premium; the package default $10^{-5}$ adds a further
safety decade at little cost.

\paragraph{Supplying the leading order.} The order $n$ is an \emph{input} of
the counting formula~\eqref{eq:count}, not a result of it, and the estimator
exists only so that a user who does not know $n$ is not blocked. Whenever $n$
is known by inspection of the model -- and it usually is -- it should simply
be passed (\code{n\_power\_max}): it is then exact rather than estimated, and
it costs nothing. Non-integer fractional orders are legitimate inputs too.

\paragraph{When the residual stops being a check.} The diagnostic has one
precondition: the \emph{convergence of the tail}. $\Zraw$ is only
\emph{expected} to be an integer when the truncated remainder
$\int_{\wmax}^{\infty}$ is negligible. For retarded systems it is, and
$\varepsilon$ behaves as described: across the retarded case studies of
Appendix~\ref{app:gallery} the median residual is $\sim10^{-4}$ and
essentially no grid point is uncertain. For \emph{neutral} systems the
integrand does not decay at all: the tail contributes an oscillation bounded
by $\arcsin(\beta)/\pi$ rather than a vanishing remainder
(Section~\ref{sec:counting}), and the residual inherits it -- on the neutral
panels of Appendix~\ref{app:gallery} the median residual is $\approx0.1$,
with the counts nevertheless correct. For that class $\varepsilon$ is
\emph{not} a usable self-check; what protects the count is the analytic
bound $\arcsin(\beta)/\pi<1/2$ instead, and enlarging $\wmax$ does not help,
because the oscillation is bounded, not decaying.

One caveat is inherited from the argument principle itself and bounds what
$\varepsilon$ can certify. A march that steps \emph{over} a peak loses
exactly $\pm\pi$ of phase, so $\Zraw$ shifts by exactly one and remains just
as close to an integer as before: $\varepsilon$ measures the
\emph{quadrature} quality and is structurally blind to a skipped peak.
Nor is this hypothetical: approaching a stability boundary the peak
half-width $|\sigma_r-\sigma|$ collapses to many orders of magnitude below
any practical step size (Table~\ref{tab:boundarystress}), and the marching
solver then strides past it and reports success. No pointwise method is
immune (Section~\ref{sec:performance-convergence}).

What rescues the situation is the second output of the same sweep. The
tracked root of Section~\ref{sec:tracking} is obtained from the minima of
$|\Dfun|^{2}$, a mechanism entirely independent of the accumulated phase,
and it remains accurate exactly where the peak is narrow -- that is, exactly
where the count is fragile. The two therefore cross-validate each other: a
skipped peak shows up as an \emph{inconsistency} between $\operatorname{sign}
(\sest-\sigma)$ and the predicate $\Zint=0$; both quantities are standard
outputs of every sweep, so the check costs one comparison.

The test must, however, be restricted to a neighbourhood of the line: far
from the boundary the tracked minimum need not belong to the \emph{dominant}
root (Section~\ref{sec:tracking}), so the tracked root may disagree with
$\Zint>0$ entirely legitimately -- applied blindly, the test fires on
\CheckNaiveRef{} of the \CheckNpoints{} points of a \emph{fully converged}
chart of the showcase system. A peak can only be stepped over where a root
lies close to the line, so the informative test is
\begin{equation}\label{eq:crosscheck}
  \operatorname{sign}(\sest-\sigma)\ \text{disagrees with}\ [\Zint=0]
  \qquad\text{and}\qquad |\sest-\sigma| < h ,
\end{equation}
with $h$ a small fraction of the spread of $\sest$ over the chart. With
$h=\CheckH$ -- well below the median $|\sest|\approx\CheckMedianSigma$ of
those spurious disagreements -- the test fires on only \CheckNearRef{} of the
same \CheckNpoints{} points. Its scope should be stated precisely. Skipping
one peak only ever changes $\Zint$ by one, so the stable/unstable
classification a chart displays is affected only where the \emph{last}
unstable root crosses -- precisely the regime in which $\sest$ is small and
the test is sharp; this is the silent error mode the residual cannot see,
and at practical tolerances it is the only one left
(Section~\ref{sec:showcase-tolerance}). What the test does not promise is to
rescue a chart run far outside the rule of thumb: at
$\mathrm{tol}=\CheckLooseTol$ the accumulated phase error itself reaches
order one, miscounts appear everywhere -- not only near the boundary -- and
the near-boundary test catches only \CheckLooseCaught{} of the
\CheckLooseFlips{} classification flips of such a run. What announces
\emph{that} failure is the residual field itself, which lights up over the
whole chart (Section~\ref{sec:showcase-tolerance}); the correct reaction is
to tighten the tolerance, not to trust any pointwise flag. The same
mechanism that flags a skipped peak can also mend it: the minimum of
$|\Dfun|^{2}$ \emph{locates} the peak precisely however narrow it is, and
Section~\ref{sec:showcase-tolerance} shows how an optional callback uses
that location to repair the count within the march itself.

\subsection{Alternative integrators within the same framework}\label{sec:alternatives}

The phase integrand~\eqref{eq:integrand} can also be fed to classical
adaptive quadrature; the package provides a Gauss--Kronrod back-end
(\pkg{QuadGK.jl}~\cite{johnson2013quadgk}) with the identical interface. For
the root count alone it is fully competitive -- within a factor of about
three of the march in evaluations in either direction
(Table~\ref{tab:evals}) -- and a reasonable choice whenever only $\Zint$ is
wanted. The decisive difference concerns the rightmost root. Adaptive
quadrature subdivides its interval recursively and visits the frequencies in
an essentially arbitrary order: there is no sweep along which a local
minimum of $|\Dfun|^{2}$ could be recognized. It is this ordered traversal
-- not accuracy, not speed, and certainly not stiffness -- that earns the
phase ODE its place as the default, because every chart in this paper is
built from the tracked root: the interpolable coloring~\eqref{eq:coloring},
the boundary objective~\eqref{eq:objective}, the refinement seeds, and the
consistency check of Section~\ref{sec:integerness} all consume $\sest$, not
$\Zint$ alone. A third back-end evaluates the integral by a fixed-step
trapezoidal rule with essentially zero overhead. It is included as a
real-time floor, but a \emph{uniform} frequency grid cannot resolve the
peak~\eqref{eq:peak} of a root close to the line, so its mean count error
converges with the step count while its \emph{worst} error over a chart does
not (Section~\ref{sec:performance-zoo}). Adaptivity is thus not a
convenience here but the very mechanism that makes the count trustworthy.